\documentclass[10pt,a4paper]{amsart}
\usepackage[margin=19mm]{geometry}
\usepackage{amsmath,amssymb,mathtools}
\usepackage[scr=rsfs,cal=euler]{mathalfa}
\usepackage{xfrac}
\usepackage[unicode,hidelinks,bookmarks=false]{hyperref}
\newcommand{\ie}{i.e.\ }
\newcommand{\cf}{cf.\ }
\newcommand{\resp}{respectively\ }
\newcommand{\Subsection}{\subsection}
\providecommand{\nobreakdash}{\nobreak\hbox{-}}
\setlength{\emergencystretch}{2em}
\multlinegap0pt
\usepackage[shortlabels]{enumitem}
\usepackage{amssymb}
\newtheorem{lemma}{Lemma}
\newtheorem{theorem}{Theorem}
\newcommand{\E}{\mathrm{e}}
\newcommand{\R}{{\mathbb R}}
\newcommand{\cF}{{\mathcal F}}
\newcommand{\cK}{{\mathcal K}}
\newcommand{\cL}{{\mathcal L}}
\newcommand{\cR}{{\mathcal R}}
\newcommand{\cU}{{\mathcal U}}
\newcommand{\ds}{\displaystyle}
\title{On Dispersive Estimates for One-Dimensional Klein--Gordon Equations}
\author{Elena Kopylova}
\date{Source: arXiv:2604.15053v1 (CC BY 4.0)}
\begin{document}
\maketitle

\noindent\emph{Source and licence.} On Dispersive Estimates for One-Dimensional Klein--Gordon Equations. Version arXiv:2604.15053v1 (CC BY 4.0), distributed by arXiv under the Creative Commons Attribution 4.0 International licence, http://creativecommons.org/licenses/by/4.0/, as recorded in the arXiv licence metadata for this exact version and shown on the abstract page https://arxiv.org/abs/2604.15053v1. Full licence notice: \texttt{licenses/arxiv-2604.15053-CC-BY-4.0.txt}.\par\smallskip

\noindent\textit{Editorial scope.} Counted unit: the article's theorem thKG3 (source0.tex lines 541--555). The article's complete original proof of it is delivered with it (source0.tex lines 557--713, one \texttt{proof} environment of 157 lines). No mathematics is written, repaired or re-derived: the statement and the proof are the article's own lines, in the article's own notation, and no retained line is substituted or re-flowed.  Retained uncounted as the source-local context the proof needs: lines 106--108; lines 290--297; lines 471--473; lines 478--480; lines 488--499.  Nothing was omitted from any retained span, so the package carries no omission record.  The retained lines cite 4 works, whose entries are transcribed verbatim from the article's own bibliography source in the e-print and delivered with short editorial labels recorded in the item's reference mapping.  No retained line contains a figure environment, an image-inclusion command or drawing code, so no image is delivered and the package declares no asset dependency.  Editorial formatting only, in addition: an AMS \texttt{amsart} preamble that restates the article's own declarations and macros used by the retained lines, namely the environments \texttt{lemma}, \texttt{theorem} on separate counters, together with the standard packages those lines need.  The retained environments print in this document as Lemma 1, Theorem 1 in order of appearance, whereas the article prints the same items with the numbering of its own full document.  The complete native source of the article as distributed in the arXiv e-print for this exact version is bundled unchanged as \texttt{sources/198577/source0.tex}.
\begin{equation}\label{V}
    |V(x)|\le C\langle x\rangle^{-\beta},\quad\quad \langle x\rangle=(1+|x|^2)^{1/2},\quad x\in\R
\end{equation}

\begin{lemma}\label{lem2} 
Let $\zeta$  be a smooth function such that $\zeta(\omega)=0$ for $\omega\le m^2+1$ and
$\zeta(\omega)=1$ for $\omega\geq m^2 + 2$. Then the decay holds
\begin{equation}\label{WD3}
\Vert  e^{-i \cK_0t} \zeta(\cK_0^2)\Vert_{{\cF}_{\sigma}\to {\cF}_{-\sigma}}
\le C (1+t)^{-\sigma},\quad\sigma>0,\quad t\ge 0.
\end{equation}
\end{lemma}

\begin{equation}\label{thh}
\| \psi(x,y, \cdot)\|_{\mathcal{A}_1} \le C,
\end{equation}

\begin{equation}\label{psi-est}
\Vert \xi(k)\partial_k\psi(x,y,k)\Vert_{\mathcal A}\le C(1+|x|)(1+|y|),
\end{equation}

\begin{equation}\label{R}
 \cR(\omega)=
   \begin{pmatrix}
 0            &  0\\
  -i          &  0
\end{pmatrix}
+
\begin{pmatrix}
 \omega &  i \\
  -i \omega^2&  \omega
\end{pmatrix} R(\omega^2-m^2).
\end{equation}

\begin{theorem}\label{thKG3}
i) Let condition \eqref{V} holds with some $\beta>2$ and let  $\zeta(x)$ be a smooth function such that $\zeta(x)=0$ for $x\le m^2+1$ and
$\zeta(x)=1$ for $x\geq m^2 + 2$.
Then  for any $\sigma>3/4$, the decay holds  
\begin{equation}\label{b1}
\big\|\E^{-it\cK}\,\zeta(\cK^2)\big\|_{{\mathcal F}_\sigma\to {\mathcal F}_{-\sigma}}
={\mathcal O}(t^{-1/2}),\quad t\to \infty.
\end{equation}
ii)
Let condition \eqref{V} holds with some $\beta>3$.  Then for any $\sigma>3/2$,
\begin{equation}\label{b2}
\big\|\E^{-it\cK}\,\zeta(\cK^2)\big\|_{{\mathcal F}_\sigma\to {\mathcal F}_{-\sigma}}
={\mathcal O}(t^{-3/2}),\quad t\to \infty.
\end{equation}
\end{theorem}

\begin{proof}
{\it Step i)}.
Substituting the  resolvent identity 
$$
R(\lambda)= R_0(\lambda)- R_0(\lambda)V R_0(\lambda)+ (R_0(\lambda)V)^2 R_0(\lambda) + (R_0(\lambda)V)^3 R(\lambda)
$$ 
into \eqref{R}, we obtain
\[
\E^{-it\cK}\,\zeta(\cK^2)=\cU_0(t)+\cU_1(t)+\cU_2(t)+{\mathcal W}(t),,
\]
where $\cU_0(t)=\E^{-it\mathbf{K_0}}\,\zeta(\cK_0^2)$, and
\begin{equation*}
\cU_j(t)=\frac 1{2\pi i}\int\limits_{\Gamma}\! \E^{-i\omega t}
([(\cR_0{\mathcal V})^j\cR](\omega+i0)-[(\cR_0{\mathcal V})^j\cR_0](\omega-i0))\zeta(\omega^2) d\omega, \quad j=1,2,
\end{equation*}
\begin{equation*}
{\mathcal W}(t)=\frac 1{2\pi i}\int\limits_{\Gamma}\! \E^{-i\omega t}
([(\cR_0{\mathcal V})^3\cR](\omega+i0)-[(\cR_0{\mathcal V})^3\cR](\omega-i0))\zeta(\omega^2) d\omega.
\end{equation*}
Here
\begin{equation}\label{cV}
{\mathcal V}(x)=\left(\begin{array}{cc}
  0               &   0
  \\
  iV(x)   &   0
  \end{array}\right).
\end{equation}
Due to Lemma \ref{lem2}, 
\begin{equation}\label{U0}
\big\| \cU_0(t)\big\|_{{\mathcal F}_{\sigma}\to {\mathcal F}_{\sigma}}
={\mathcal O}(t^{-\sigma}),\quad t\to \infty,\quad \sigma>0.
\end{equation}
Let us prove that for $\beta>2$ %and $\sigma=\beta/2$
\begin{equation}\label{U1}
\big\| \cU_j(t)\big\|_{{\mathcal F}_{\beta/2}\to {\mathcal F}_{-\beta/2}}
={\mathcal O}(t^{-\beta/2}),\quad t\to \infty,\quad j=1,2.
\end{equation}
\smallskip\\
For arbitrary  $\Psi_0\in {\mathcal F}_{\beta/2}$, denote  $\Psi_{j+1}(t)=\cU_j(t)\Psi_0$, $j=0,1,2$.
%%%%%%%%%%%%%%%%%%
\begin{lemma}(cf. \cite[Lemma 3.8]{KK10}, \cite[Lemma 36.6]{KK12}). 
The convolution representation hold
\begin{equation}\label{conv}
\Psi_{j+1}(t)=
 i\int_0^t \cU_0(t-\tau){\mathcal V} \Psi_j(\tau)~d\tau,~~~~t\in\R,\quad j=1,2,
\end{equation}
where the integral converges in ${\mathcal F}_{-\beta/2}$.
\end{lemma}
By Lemma \ref{lem2},
$\Vert\Psi_1\Vert_{\cF_{-\beta/2}}\le C (1+|t|)^{-\beta/2}\Vert \Psi_0\Vert_{\cF_{\beta/2}}$.
 Applying  Lemma \ref{lem2} to the integrand in (\ref{conv}), 
 we obtain for $\sigma=\beta/2$
 \begin{equation*}
 \Vert \cU_0(t-\tau){\mathcal V} \Psi_1(\tau)\Vert_{\cF_{-\sigma}}\le 
 \frac{C\Vert {\mathcal V} \Psi_1(\tau)\Vert_{\cF_{\sigma}}}{(1\!+|t-\!\tau|)^{\sigma}}
 \le\frac{C_1\Vert \Psi_1(\tau)\Vert_{\cF_{-\sigma}}}{(1\!+|t-\!\tau|)^{\sigma}}
  \le\!\frac{C_2\Vert \Psi_0\Vert_{\cF_{\sigma}}}{(1\!+|t-\!\tau|)^{\sigma}(1\!+|\tau|)^{\sigma}}
\end{equation*}  
Therefore, integrating here in $\tau$, we obtain 
\begin{equation*}
  \Vert \Psi_2(t)\Vert_{\cF_{-\beta/2}} \le C (1+|t|)^{-\beta/2}\Vert \Psi_0\Vert_{\cF_{\beta/2}}.
\end{equation*}
Similarly,
\begin{equation*}
  \Vert \Psi_3(t)\Vert_{\cF_{-\beta/2}} \le C (1+|t|)^{-\beta/2}\Vert \Psi_0\Vert_{\cF_{\beta/2}}.
\end{equation*}
Hence, \eqref{U1} follows.
\smallskip\\
It remains to estimate the operator  ${\mathcal W}(t)$ with the kernel
\begin{equation}\label{W}
\!\!\!\!\!{\mathcal W}(t,x,y)
=\!\int\limits_{\R^3} \frac{{\mathbf V}(u,\rm w,z)}{16 \pi i}\left(\int\limits_\R\! \xi(k){\mathcal M}_t(k)
\frac{\E^{ikp(x,y,u,\rm w,z)}}{k^3}\psi(y,z,k)dk\right) dud{\rm w} dz,
\end{equation}
where we denote  $p(x,y,u,\rm w,z)=|x-u|+|u-{\rm w}|+|{\rm w}-z|+|z-y|$,  $\xi(k)=\zeta(k^2+m^2)$,
${\mathbf V}(u,\rm w,z)=V(u)V(\rm w)V(z)$.
We will apply the following version of \cite[Lemma 2]{MSW}:
%%%%%%%%%%%%%%%%%%%%%%%%%%%%%%%%%%%%%%%%%%%%%%%%%%%%%%%%%%%%%%%%%%%%%%%%%%%%%%%%%
%%%%%%%%%%%%%%%%%%%%%%%%%%%%%%%%%%%%%%%%%%%%%%%%%%%%%%%%%%%%%%%%%%%%%%%%%%%%%%%%%
\begin{lemma}\label{l1} (cf. \cite[Lemma 5.5]{EKTM})
Let $\eta(k)$, $k\geq 1$, be a smooth function such that
 $|\eta^{(j)}(k)|\le k^{-j}$ for $j=0,1$. Then for any $g(k)\in \mathcal A_1$,
 $\alpha> 3/2$ and $t\ge 1$
\begin{equation}\label{OI-est}
\sup\limits_{p\in\R}\Big|\int_{1}^\infty\eta(k)\frac{\E^{\pm it\omega(k)+ ik p}}
{k^\alpha}g(k)dk\Big|\le C t^{-1/2}\Vert g\Vert_{\mathcal A_1} ,\quad \omega(k)=\sqrt{k^2+m^2}.
\end{equation}
\end{lemma}
Applying this lemma with  $g(k)=\psi(y,z,k)$, $p= p(x,u,\rm w,z,y)$, and taking into account \eqref{thh}, we get
\begin{equation}\label{born11}
\sup\limits_{x,y}|{\mathcal W}^{ij}(t,x,y)|\le C t^{-1/2}\|V\|_{L^1}^3\le C_1t^{-1/2},\quad t\ge 1,\quad i,j=1,2.
\end{equation}
Here ${\mathcal W}^{ij}$ are the entries of the matrix ${\mathcal W}$. Further,
\begin{eqnarray}\nonumber
\!\!\!\!\!\!\!\!\!\!\!\! &&\!\!\!\!\!\!\!\!\partial_x{\mathcal W}(t,x,y)\\
\nonumber
\!\!\!\!\!\!\!\!\!\!\!\!&&\!\!\!\!\!\!\!\!=\int\limits_{\R^3} {\rm sgn} (x\!-\!u)\frac{{\mathbf V}(u,\rm w,z)}{16 \pi}\left(\int\limits_{\R} \xi(k)
{\mathcal M}_t(k)
\frac{\E^{ikp(x,y,u,\rm w,z)}}{k^2}\psi(y,z,k)dk\right) dud{\rm w} dz.
\end{eqnarray}
Hence, similarly to \eqref{born11},
\begin{equation}\label{born12}
\sup\limits_{x,y}|\partial_x{\mathcal W}^{1j}(t,x,y)|\le C t^{-1/2},\quad t\ge 1,\quad j=1,2.
\end{equation}
Bounds \eqref{born11} and \eqref{born12} imply that
\[%\begin{equation}\label{U2}
\Vert {\mathcal W}(t)\Vert_{\cL^{!}\to \cL^\infty}\le C t^{-1/2}, \quad t\ge 1.
\]%\end{equation}
Together with  \eqref{U0} and \eqref{U1} this gives \eqref{b1}.
\smallskip\\
{\it Step ii)}.
Now we rewrite \eqref{W} as follows
\begin{eqnarray}\nonumber
\!\!\!\!\!\!\!\!\!\!&&\!\!\!\!\!\!\!\!\!\!\!\!\!{\mathcal W}(t,x,y)\\
\label{W1}
\!\!\!\!\!\!\!\!\!\!&&\!\!\!\!\!\!\!\!\!\!\!\!\!
= \!\sum_{\pm} \!\int\limits_{\R^3} \frac{{\mathbf V}(u,{\rm w},z)}{16 \pi i}\left(\int\limits_\R\! \xi(k)A_{\pm}(k)
\frac{\E^{it\omega(k)+ikp(x,y,u,{\rm w},z)}}{k^3}\, \psi(y,z,k)dk\!\right)\! dud{\rm w} dz,
\end{eqnarray}
where
\[
A_{\pm}(k)=\begin{pmatrix}
1  & \mp\ds\frac{i}{\omega(k)}\\
  \pm i\omega(k)& 1
\end{pmatrix},
\]
Integrating by parts in inner integral of \eqref{W1}, we get
\begin{eqnarray}\nonumber
&&F_{\pm}(t,x,y,u,{\rm w},z):=\int \xi(k)A_{\pm}(k)\frac{\E^{it\omega(k)+ikp(x,y,u,{\rm w},z)}}{k^3}\, \psi(y,z,k)dk\\
\nonumber
&&=\frac{i}{t}\int \E^{it\omega(k)}\frac{\partial}{\partial k}
\Big[\E^{ikp(x,y,u,{\rm w},z)}\xi(k)\frac{\omega(k)}{k^4}\,A_{\pm}(k) \psi(y,z,k)\Big]dk.
\end{eqnarray}
Recall that in the case $V\in L^1_2$,
$\Vert \frac{\partial}{\partial k}\psi(z,y,k)\Vert_{\mathcal A}\le C(1+|z|)(1+|y|)$  by \eqref{psi-est}. 
Moreover,
\[
|x-u|+|u-{\rm w}|+|{\rm w}-z|+|z-y|\le (1+|x|)(1+|y|)(1+2|z|) (1+2|u|)(1+2|{\rm w}|).
\] 
Hence, by Lemma \ref{l1}, 
\[
|F_{\pm}(t,x,y,u,{\rm w},z) |\le Ct^{-3/2} (1+|x|)(1+|y|)(1+|u|)(1+|{\rm w}|)(1+|z|),\quad t\ge 1.
\]
The last estimate and \eqref{W1}  then imply
\[
|{\mathcal W}^{ij}(x,y)|\le  Ct^{-3/2} (1+|x|)(1+|y|),\quad t\ge 1,\quad i,j=1,2.
\] 
Similarly,
\[
|\partial_x {\mathcal W}^{1j}(x,y)|\le  Ct^{-3/2} (1+|x|)(1+|y|),\quad t\ge 1,\quad j=1,2.
\]
Therefore,
\[
\Vert {\mathcal W}(t)\Vert_{\cL^1_1\to\cL^\infty_1}\le  Ct^{-3/2}, \quad t\ge 1.
\]
Together with  \eqref{U0} and \eqref{U1} this gives \eqref{b2}.
\end{proof}

\begin{thebibliography}{99}
\bibitem[EKTM]{EKTM} I. Egorova, E. Kopylova, V. A. Marchenko and G. Teschl, On the dispersion estimates for the one-dimensional Schrdinger and Klein-Gordon equa- tions. Revizited. {\em Russian Math. Surveys} {\bf 71} (2016), no. (429), 391-415.
\bibitem[KK10]{KK10} A. Komech and E. Kopylova, Weighted energy decay for 3D Klein-Gordon equation, {\em J. of Diff. Eqns.}, {\bf 248} (2010), no. 3, 501--520.
\bibitem[KK12]{KK12} A. Komech and E. Kopylova, Dispersion decay and scattering theory, John Willey \& Sons, Hoboken, New Jersey, 2012.
\bibitem[MSW]{MSW} B. Marshall, W. Strauss, and S. Wainger, $L^p-L^q$ estimates for the Klein--Gordon equation, {\em J. Math. Pures et Appl.} {\bf 59} (1980), 417--440.
\end{thebibliography}
\end{document}
